% =====================================================================
%  ch01b.tex —— 1.2 导出范畴与导出函子；1.3 概形上层的导出范畴
%                （原书 p.523 中“1.2.”标题起，至 p.530 末、第 2 章标题前）
%  底本：D. O. Orlov, Russian Math. Surveys 58:3 (2003) 511--591
% =====================================================================
% [原文 p.523]
\section{导出范畴与导出函子}
\takeaway{本节先在加性范畴上造出复形范畴、同伦范畴与映射锥，证明 $\mathbf{H}(\cA)$ 以标准三角为特异三角构成三角范畴；再以零调复形作局部化得 $\mathbf{D}(\cA)=\mathbf{H}(\cA)/\cN$，并用内射消解与右适应子范畴给出导出函子 $\mathbf{R}F$、$\mathbf{L}F$ 有定义的条件。}

% [Tutor Note: 原书以粗体 $\mathbf{C}$、$\mathbf{H}$、$\mathbf{D}$ 区分复形范畴、同伦范畴、导出范畴，与花体 $\cC$、$\cH$ 不同；本译本照此，一律排作 $\mathbf{C}(\cA)$、$\mathbf{H}(\cA)$、$\mathbf{D}(\cA)$。]
设 $\cA$ 是加性范畴。我们用 $\mathbf{C}(\cA)$ 记微分复形的范畴。它的对象是复形
\begin{equation*}
M^{\cdot}=\bigl(\cdots\longrightarrow M^{p}\xrightarrow{\ d^{p}\ }M^{p+1}\longrightarrow\cdots\bigr)
\end{equation*}
其中 $M^{p}\in\cA$ 对 $p\in\Z$，且 $d^{2}=0$；而态射 $f:M^{\cdot}\to N^{\cdot}$ 是一族 $\cA$ 中的态射 $f^{p}:M^{p}\to N^{p}$，它们与微分交换，即
\begin{equation*}
d_{N}f^{p}-f^{p+1}d_{M}=0
\end{equation*}
对任意 $p$ 成立。

我们用 $\mathbf{C}^{+}(\cA)$、$\mathbf{C}^{-}(\cA)$、$\mathbf{C}^{b}(\cA)$ 记 $\mathbf{C}(\cA)$ 的如下全子范畴，它们分别由满足下述条件的复形 $M^{\cdot}$ 构成：对所有 $p\ll 0$ 有 $M^{p}=0$；对所有 $p\gg 0$ 有 $M^{p}=0$；对所有 $p\gg 0$ 且所有 $p\ll 0$ 有 $M^{p}=0$。

我们说复形的态射 $f:M^{\cdot}\to N^{\cdot}$ 是\emph{零伦的}（null-homotopic），如果存在一族态射 $h^{p}:M^{p+1}\to N^{p}$，使得对所有 $p\in\Z$ 有 $f^{p}=d_{N}h^{p}+h^{p+1}d_{M}$。我们定义\emph{同伦范畴} $\mathbf{H}(\cA)$ 为与 $\mathbf{C}(\cA)$ 有相同对象的范畴，而 $\mathbf{H}(\cA)$ 中的态射是复形之间态射 $f$ 的类 $\overline{f}$，模去零伦态射。

我们由规则
\begin{equation*}
(M[1])^{p}=M^{p+1},\qquad d_{M[1]}=-d_{M}
\end{equation*}
定义平移函子 $[1]:\mathbf{H}(\cA)\to\mathbf{H}(\cA)$。

我们把 $\mathbf{H}(\cA)$ 中的\emph{标准三角}定义为序列
\begin{equation*}
L\xrightarrow{\ \overline{f}\ }M\xrightarrow{\ \overline{g}\ }Cf\xrightarrow{\ \overline{h}\ }L[1],
\end{equation*}
其中 $f:L\to M$ 是某个复形态射，$Cf=M\oplus L[1]$ 是 $\mathbf{C}(\cA)$ 的分次对象，其微分为
\begin{equation*}
d_{Cf}=\begin{pmatrix} d_{M} & f\\ 0 & -d_{L}\end{pmatrix},
\end{equation*}
$g$ 是典范嵌入 $M\to Cf$，$-h$ 是典范投影。照例，$Cf$ 称为 $f$ 的\emph{映射锥}。

\yizhu{映射锥 $Cf=M\oplus L[1]$ 的微分矩阵里，$f$ 那一项正是把 $f$ 与 $L$ 的平移耦合起来的地方。由三元组的长正合列可见，$H^{n}(Cf)$ 夹在 $H^{n}(L)$ 与 $H^{n}(M)$ 之间并与 $H^{n}(f)$ 相衔接，故 $f$ 是拟同构当且仅当 $Cf$ 零调——这正是下文定义拟同构的方式。三个箭头都取同伦类并带上横线，说明标准三角只在 $\mathbf{H}(\cA)$ 中有意义，而非 $\mathbf{C}(\cA)$ 中的序列。【来源：原书 pp.523--524；Verdier, Astérisque 239 (1996)。】}

% [原文 p.524]
\begin{lemma}{1.2.1}
以 $[1]$ 为平移函子、以与标准三角同构的三角构成的类为特异三角，范畴 $\mathbf{H}(\cA)$ 是三角范畴。
\end{lemma}

我们用 $\mathbf{H}^{+}(\cA)$、$\mathbf{H}^{-}(\cA)$、$\mathbf{H}^{b}(\cA)$ 记 $\mathbf{C}^{+}(\cA)$、$\mathbf{C}^{-}(\cA)$、$\mathbf{C}^{b}(\cA)$ 在 $\mathbf{H}(\cA)$ 中的像。这些范畴也是三角范畴。

现在设 $\cA$ 是阿贝尔范畴。为定义阿贝尔范畴的导出范畴，须回忆零调复形与拟同构这两个概念。对任一复形 $N^{\cdot}$ 与每个 $p\in\Z$，上同调 $H^{p}(N^{\cdot})\in\cA$ 定义为 $\Ker d^{p}/\im d^{p-1}$。于是，对任一整数 $p$ 都有加性函子 $H^{p}:\mathbf{C}(\cA)\to\cA$，它把复形 $N^{\cdot}$ 送到它的第 $p$ 个上同调 $H^{p}(N^{\cdot})\in\cA$。

我们说复形 $N^{\cdot}\in\mathbf{C}(\cA)$ 在第 $n$ 项处\emph{零调}（acyclic），如果 $H^{n}(N^{\cdot})=0$；若它的上同调全部消没，即对 $n\in\Z$ 有 $H^{n}(N^{\cdot})=0$，就简称它为\emph{零调的}。我们用 $\cN$ 记 $\mathbf{H}(\cA)$ 中由所有零调复形构成的全子范畴。下面的引理几乎是显然的。

\begin{lemma}{1.2.2}
子范畴 $\cN$ 是 $\mathbf{H}(\cA)$ 的全三角子范畴。
\end{lemma}

我们说 $\mathbf{H}(\cA)$ 中的态射 $f:X\to Y$ 是\emph{拟同构}，如果它的映射锥是零调复形。换言之，$f$ 是拟同构当且仅当它在上同调上诱导的映射是同构。设 $\mathrm{Quis}$ 为与 $\cN$ 相伴的乘法系，即由所有拟同构构成的系。

\begin{definition}{1.2.3}
阿贝尔范畴 $\cA$ 的\emph{导出范畴} $\mathbf{D}(\cA)$ 定义为同伦范畴 $\mathbf{H}(\cA)$ 关于所有零调复形构成的子范畴的局部化，即
\begin{equation*}
\mathbf{D}(\cA):=\mathbf{H}(\cA)/\cN=\mathbf{H}(\cA)[\mathrm{Quis}^{-1}].
\end{equation*}
\end{definition}

对 $\ast\in\{+,-,b\}$，我们用同样的方式定义相应的导出范畴 $\mathbf{D}^{\ast}(\cA)$，即局部化 $\mathbf{H}^{\ast}(\cA)/(\mathbf{H}^{\ast}(\cA)\cap\cN)$。

\begin{lemma}{1.2.4}
对 $\ast\in\{+,-,b\}$，典范函子
\begin{equation*}
\mathbf{D}^{\ast}(\cA)\longrightarrow\mathbf{D}(\cA)
\end{equation*}
与 $\mathbf{D}(\cA)$ 中由那些分别在 $n\ll 0$ 时、在 $n\gg 0$ 时、以及在 $n\ll 0$ 与 $n\gg 0$ 时零调的复形构成的全子范畴给出等价。子范畴 $\mathbf{H}^{+}(\cA)$ 关于拟同构类在 $\mathbf{H}(\cA)$ 中是右共尾的，而 $\mathbf{H}^{-}(\cA)$ 是左共尾的。
\end{lemma}

\yizhu{局部化 $\mathbf{H}(\cA)[\mathrm{Quis}^{-1}]$ 的含义是：把拟同构形式地变成同构，即为每个拟同构添一个逆，再对由这些逆与原态射拼出的三角关系取商。这一步之所以可行，靠的是 $\mathrm{Quis}$ 满足乘法系的 Ore 条件（1.1 节已列出局部化公理 L1--L4），于是分式可写成 $s^{-1}f$ 而不必担心分母。$\mathbf{D}^{+}$、$\mathbf{D}^{-}$、$\mathbf{D}^{b}$ 并不是这样局部化出来的商，而是 $\mathbf{D}(\cA)$ 的全子范畴，其“足够大”由引理 1.2.4 的共尾性保证。【来源：原书 pp.524--525；Gabriel--Zisman, \emph{Calculus of Fractions and Homotopy Theory}, Springer 1967。】}

设阿贝尔范畴 $\cA$ 有足够多的内射对象，即每个对象都嵌入某个内射对象。我们用 $\cI$ 记 $\cA$ 中由内射对象构成的全子范畴。此时可以证明，复合函子
\begin{equation*}
\mathbf{H}^{+}(\cI)\longrightarrow\mathbf{H}^{+}(\cA)\xrightarrow{\ Q\ }\mathbf{D}^{+}(\cA)
\end{equation*}
是范畴等价（见 [17]）。类似地，若阿贝尔范畴 $\cA$ 有足够多的投射对象，则复合函子
\begin{equation*}
\mathbf{H}^{-}(\cP)\longrightarrow\mathbf{H}^{-}(\cA)\xrightarrow{\ Q\ }\mathbf{D}^{-}(\cA)
\end{equation*}
是等价，其中 $\cP$ 是 $\cA$ 中由投射对象构成的全子范畴。

\yizhu{这两条等价是“用内射消解计算右导出函子”的全部根据：$\mathbf{H}^{+}(\cI)$ 中的对象已是内射复形，嵌入 $\mathbf{H}^{+}(\cI)\to\mathbf{H}^{+}(\cA)$ 是全忠实的，而 $Q$ 把拟同构变成同构；两者复合为等价，等于说 $\mathbf{D}^{+}(\cA)$ 的每个对象都有内射复形代表。于是 $\mathbf{R}F$ 只须在 $\mathbf{H}^{+}(\cI)$ 上按 $F$ 直接作用即可，结果的唯一性由拟同构变成同构保证。有足够多投射对象时对偶地得到左导出函子。【来源：原书 p.524 及其所引 [17]；按所引章节编号（I.4.6、II.5.6 等）判断，[17] 当为 Hartshorne, \emph{Residues and Duality}, LNM 20, 1966。】}

% [原文 p.525]
设 $F:\cA\to\cB$ 是阿贝尔范畴之间的加性函子（不必正合）。于是 $F$ 以显然的方式诱导正合函子 $\mathbf{H}(\cA)\to\mathbf{H}(\cB)$，我们仍用同一符号 $F$ 记它。上一节所述的（右）导出函子的一般构造给出一个函子 $\mathbf{R}F$，它定义在 $\mathbf{D}(\cA)$ 的某个全三角子范畴上，取值于 $\mathbf{D}(\cB)$。对左导出函子也有同样的说法。我们把 $F$ 的第 $n$ 个右导出函子（相应地，左导出函子）定义为上同调
\begin{equation*}
\mathbf{R}^{n}FX=H^{n}(\mathbf{R}FX)\qquad\bigl(\text{相应地，}\ \mathbf{L}_{n}FX=H^{-n}(\mathbf{L}FX)\bigr),
\end{equation*}
其中 $n\in\Z$。

在应用中，右伴随函子通常恰好在子范畴 $\mathbf{D}^{+}(\cA)$ 上有定义。利用引理 1.1.10 与引理 1.2.4，可以说：函子 $\mathbf{R}F$ 在 $\mathbf{D}^{+}(\cA)$ 上的限制，与 $F$ 在 $\mathbf{H}^{+}(\cA)\subset\mathbf{H}(\cA)$ 上的限制的导出函子一致。

下面说明右导出函子 $\mathbf{R}F$ 在整个范畴 $\mathbf{D}^{+}(\cA)$ 上有定义的条件。我们说全加性子范畴 $\cR\subset\cA$ 关于函子 $F$ 是\emph{右适应的}，如果
\begin{enumerate}[label=\alph*),leftmargin=2.4em,itemsep=1pt,topsep=2pt]
\item $F$ 把 $\mathbf{C}^{+}(\cR)$ 中的零调复形送到零调复形；
\item $\cA$ 的每个对象都嵌入 $\cR$ 中某个对象。
\end{enumerate}

我们说 $\cR$ 的对象是\emph{右 $F$-零调的}。若存在关于 $F$ 右适应的子范畴 $\cR$，通常就说 $\cA$ 有足够多的（右）$F$-零调对象。

\yizhu{“适应子范畴”是内射对象这一特例的公理化：条件 b) 保证每个对象都有 $\cR$-消解，条件 a) 保证这套消解在 $F$ 之下仍然零调，于是把 $F$ 作用在消解上就得到与消解选取无关的结果。内射对象关于任意加性函子都右适应，因为内射层在左正合函子下的像是零调的；所以“内射消解”只是 $\cR=\cI$ 的特例。概形上 $\cI$ 常常跑出 $\Qcoh(X)$，这时适应子范畴就取代了它。【来源：原书 p.525；同 [17]。】}

设 $F:\cA\to\cB$ 是存在右适应子范畴 $\cR\subset\cA$ 的函子。应用引理 1.1.12 容易验证，$\cR$ 中对象构成的每个右有界复形都是右 $F$-分裂的。由条件 b) 推出，对每个对象 $X\in\mathbf{H}^{+}(\cA)$ 都存在拟同构 $X\to X'$，其中 $X'\in\mathbf{H}^{+}(\cR)$。作为推论，我们看到典范函子
\begin{equation*}
\mathbf{H}^{+}(\cR)[\mathrm{Quis}^{-1}]\longrightarrow\mathbf{D}^{+}(\cA)
\end{equation*}
是三角范畴的等价。

\begin{lemma}{1.2.5}
设 $F$ 是存在右适应子范畴 $\cR\subset\cA$ 的函子。则函子 $\mathbf{R}F$ 在 $\mathbf{D}^{+}(\cA)$ 上有定义，且对任一左有界复形 $X$ 都有同构 $\mathbf{R}FX\xrightarrow{\ \sim\ }X'$，其中 $X\to X'$ 是拟同构，$X'\in\mathbf{H}^{+}(\cA)$。
\end{lemma}
% [Tutor Note: 引理 1.2.5 的页图与文字层一致，均印作 $\mathbf{R}FX\xrightarrow{\sim}X'$ 与 $X'\in\mathbf{H}^{+}(\cA)$。按上下文（$X\to X'$ 为拟同构，且须对 $X'$ 直接施用 $F$ 才有意义）此处当读作 $\mathbf{R}FX\xrightarrow{\sim}FX'$、$X'\in\mathbf{H}^{+}(\cR)$。本译本照录原文，仅注明于此。]
% [Tutor Note: 原书 left bounded / right bounded 的用法为：left bounded 指 $p\ll 0$ 时消没（即有界在下、属 $\mathbf{D}^{+}$），right bounded 指 $p\gg 0$ 时消没（即有界在上、属 $\mathbf{D}^{-}$），与部分文献的习惯相反。本译本照原文直译为“左有界”“右有界”，读者宜以消没方向为准。]

若 $\cA$ 有足够多的内射对象，则由所有内射对象构成的全子范畴 $\cI\subset\cA$ 关于任一加性函子都是右适应的。此时只要把 $F$ 作用于复形 $X$ 的内射消解 $X'$，就能算出右导出函子 $\mathbf{R}FX$。

对偶地，可以引入关于函子 $F$ 左适应的子范畴这一概念。若这样的子范畴存在，则左导出函子 $\mathbf{L}F:\mathbf{D}^{-}(\cA)\to\mathbf{D}(\cB)$ 有定义。

% [原文 p.525]
\section{概形上层的导出范畴}
\takeaway{本节把 1.2 的机器用到概形上：比较 $\cO_X$-Mod、$\Qcoh(X)$、$\coh(X)$ 三层的导出范畴，说明为何要在拟凝聚层里做内射消解；随后定义 $\mathbf{L}f^{*}$、$\otL$、$\mathbf{R}f_{*}$、$\mathbf{R}\mathcal{H}om$，并列出投影公式、平坦基变换与式 (5) 三条运算律。}

任一概形都可以配上若干个层的阿贝尔范畴。设 $X$ 是域 $k$ 上的概形，结构层为 $\cO_X$。我们用 $\cO_X$-Mod 记 Zariski 拓扑中所有 $\cO_X$-模层构成的阿贝尔范畴。范畴 $\cO_X$-Mod 具有所有极限与
% [原文 p.526]
余极限，并有一个生成元集。正向余极限正合。因此，范畴 $\cO_X$-Mod 是 Grothendieck 阿贝尔范畴，有足够多的内射对象（见 [15]、[42]，Exp. IV）。

从现在起，我们只考虑诺特概形（不过下文处理的许多事实在更一般的情形下也成立）。我们用 $\Qcoh(X)$ 记 $\cO_X$-Mod 中由拟凝聚层构成的全阿贝尔子范畴。在诺特概形 $X$ 上，每个拟凝聚层都是其有限型子层的正向余极限（见 [16]，EGA1，9.4）。此时范畴 $\Qcoh(X)$ 有一个生成元集，是 Grothendieck 阿贝尔范畴，因而有足够多的内射对象。

我们能给概形 $X$ 配上的第三个范畴是凝聚层范畴 $\coh(X)$；它是 $\Qcoh(X)$ 的全阿贝尔子范畴。虽然（拟）凝聚层的定义是局部的，其实它们并不依赖拓扑。例如，我们不必只考虑 Zariski 拓扑，也可以考虑 étale 拓扑或平坦拓扑。此时，尽管 $\cO_X$-模层的概念依赖拓扑的选取，（拟）凝聚层却不依赖（见 [36]）。特别地，对仿射概形 $X$，范畴 $\Qcoh(X)$ 等价于与 $X$ 相应的代数上的模范畴。
% [Tutor Note: 原书此处印作 “etale”，无重音符；本译本按通行拼写作 étale，并在此注明。]

\yizhu{三层的关系是 $\coh(X)\subset\Qcoh(X)\subset\cO_X\text{-}\Mod$。最外层是 Grothendieck 范畴，内射对象足够，但凝聚层在其中做内射消解会跑出 $\coh(X)$，所以 $\coh(X)$ 自身没有足够多内射对象；中间的 $\Qcoh(X)$ 在诺特概形上仍有足够多内射对象，又对 $\coh(X)$ 足够“大”，于是成为计算凝聚层导出函子的工作台。命题 1.3.1--1.3.3 正是在这三层之间搬运全忠实性。【来源：原书 p.526；Hartshorne, \emph{Algebraic Geometry}, GTM 52, II.5。】}

下文我们将着眼于凝聚层范畴，更确切地说，着眼于凝聚层的导出范畴。然而，由于 $\coh(X)$ 没有足够多的内射对象，在构造导出函子时我们要借助范畴 $\Qcoh(X)$ 与 $\cO_X$-Mod。

对诺特概形 $X$，阿贝尔范畴的全嵌入 $\Qcoh(X)\hookrightarrow\cO_X$-Mod 把内射对象送到内射对象。由此可用一个简单步骤（见 [17]，I.4.6；[43]，附录 B）推出，三角子范畴 $\mathbf{H}^{+}(\Qcoh)$ 在三角范畴 $\mathbf{H}^{+}(\cO_X\text{-}\Mod)$ 中是右共尾的。于是应用引理 1.1.6，得到如下断言。

\begin{proposition}{1.3.1}
（[17]，[41]，Exp. II）若 $X$ 是诺特概形，则典范函子
\begin{equation*}
\mathbf{D}^{+}(\Qcoh(X))\longrightarrow\mathbf{D}^{+}(\cO_X\text{-}\Mod)
\end{equation*}
是全忠实的，并与由具有拟凝聚上同调的复形构成的全子范畴 $\mathbf{D}^{+}(\cO_X\text{-}\Mod)_{\Qcoh}\subset\mathbf{D}^{+}(\cO_X\text{-}\Mod)$ 给出等价。
\end{proposition}

在概形上添加额外条件后，我们也能证明无界导出范畴的类似断言。

\begin{proposition}{1.3.2}
（[41]，Exp. II）若 $X$ 是有限维诺特概形，则典范函子
\begin{equation*}
\mathbf{D}(\Qcoh(X))\longrightarrow\mathbf{D}(\cO_X\text{-}\Mod)
\end{equation*}
是全忠实的，并与由具有拟凝聚上同调的复形构成的全子范畴 $\mathbf{D}(\cO_X\text{-}\Mod)_{\Qcoh}\subset\mathbf{D}(\cO_X\text{-}\Mod)$ 给出等价。
\end{proposition}

% [原文 p.527]
证明用到如下事实：嵌入函子有右伴随 $Q:\cO_X\text{-}\Mod\to\Qcoh(X)$，而对有限维概形，这个函子有有限的上同调维数（见 [41]，II.3.7）。

现在考虑阿贝尔范畴的嵌入 $\coh(X)\subset\Qcoh(X)$。对导出范畴之间的典范函子，也有与上述类似的已知断言；不过这些断言只涉及右有界导出范畴。

\begin{proposition}{1.3.3}
（[41]，Exp. II）对诺特概形 $X$，典范函子
\begin{equation*}
\mathbf{D}^{-}(\coh(X))\longrightarrow\mathbf{D}^{-}(\Qcoh(X))
\end{equation*}
是全忠实的，并与全子范畴 $\mathbf{D}^{-}(\Qcoh(X))_{\coh}$ 给出等价。
\end{proposition}

把这个命题与命题 1.3.1、1.3.2 结合，得到如下推论。

\begin{corollary}{1.3.4}
（[41]，Exp. II）设 $X$ 是诺特概形（相应地，有限维诺特概形）。则典范函子
\begin{equation*}
\mathbf{D}^{b}(\coh(X))\longrightarrow\mathbf{D}^{b}(\cO_X\text{-}\Mod)\qquad
\bigl(\text{相应地，}\ \mathbf{D}^{-}(\coh(X))\longrightarrow\mathbf{D}^{-}(\cO_X\text{-}\Mod)\bigr)
\end{equation*}
是全忠实的，并与全子范畴 $\mathbf{D}^{b}(\cO_X\text{-}\Mod)_{\coh}$（相应地，$\mathbf{D}^{-}(\cO_X\text{-}\Mod)_{\coh}$）给出等价。
\end{corollary}

下面描述概形上层的导出范畴之间的主要导出函子。设 $f:X\to Y$ 是诺特概形的态射。存在逆像函子
\begin{equation*}
f^{*}:\cO_Y\text{-}\Mod\longrightarrow\cO_X\text{-}\Mod,
\end{equation*}
它是右正合的。由于 $\cO_Y$-Mod 有足够多的平坦 $\cO_Y$-模，而它们是 $f^{*}$-零调的，可知左导出函子
\begin{equation*}
\mathbf{L}f^{*}:\mathbf{D}^{-}(\cO_Y\text{-}\Mod)\longrightarrow\mathbf{D}^{-}(\cO_X\text{-}\Mod)
\end{equation*}
有定义。容易验证，$\mathbf{L}f^{*}$ 把范畴 $\mathbf{D}^{-}(\cO_Y\text{-}\Mod)_{\Qcoh}$ 与 $\mathbf{D}^{-}(\cO_Y\text{-}\Mod)_{\coh}$ 分别送到 $\mathbf{D}^{-}(\cO_X\text{-}\Mod)_{\Qcoh}$ 与 $\mathbf{D}^{-}(\cO_X\text{-}\Mod)_{\coh}$。于是，对有限维诺特概形，我们在拟凝聚层与凝聚层的右有界导出范畴上得到导出逆像函子 $\mathbf{L}f^{*}$。

若 $f^{*}$ 有有限的上同调维数（此时我们说 $f$ 有有限 Tor 维数），就可以把导出函子 $\mathbf{L}f^{*}$ 延拓到无界导出范畴。此外，若 $f$ 有有限 Tor 维数，则导出逆像函子把有界导出范畴送到有界导出范畴。特别地，我们有函子
\begin{equation*}
\mathbf{L}f^{*}:\mathbf{D}^{b}(\cO_Y\text{-}\Mod)_{\coh}\longrightarrow\mathbf{D}^{b}(\cO_X\text{-}\Mod)_{\coh}.
\end{equation*}

% [原文 p.528]
设 $\cE,\cF\in\mathbf{C}(\cO_X\text{-}\Mod)$ 是 $\cO_X$-模的两个复形。我们把张量积 $\cE\otimes\cF$ 定义为与二重复形 $\cE^{p}\otimes\cF^{q}$ 相伴的复形，即
\begin{equation*}
(\cE\otimes\cF)^{n}=\sum_{p+q=n}\cE^{p}\otimes\cF^{q},
\end{equation*}
其微分为 $d=d_{\cE}+(-1)^{n}d_{\cF}$。复形态射之间的同伦可以延拓到张量积上，于是得到函子
\begin{equation*}
\cE\otimes(\,\cdot\,):\mathbf{H}(\cO_X\text{-}\Mod)\longrightarrow\mathbf{H}(\cO_X\text{-}\Mod).
\end{equation*}

现在设 $\cE\in\mathbf{C}^{-}(\cO_X\text{-}\Mod)$。范畴 $\mathbf{H}^{-}(\cO_X\text{-}\Mod)$ 有足够多的关于函子 $\cE\otimes(\,\cdot\,)$ 左分裂的对象；事实上，有界在上（右有界）的平坦 $\cO_X$-模复形就有这个性质。因此存在左导出函子
\begin{equation*}
\cE\otL:\mathbf{D}^{-}(\cO_X\text{-}\Mod)\longrightarrow\mathbf{D}^{-}(\cO_X\text{-}\Mod).
\end{equation*}

若 $\cE_{1}$ 与 $\cE_{2}$ 拟同构，则 $\cE_{1}\otL$ 与 $\cE_{2}\otL$ 同构。事实上我们得到二变元函子
\begin{equation*}
\otL:\mathbf{D}^{-}(\cO_X\text{-}\Mod)\times\mathbf{D}^{-}(\cO_X\text{-}\Mod)\longrightarrow\mathbf{D}^{-}(\cO_X\text{-}\Mod),
\end{equation*}
它对两个变元都正合。张量积的导出函子显然是结合且对称的。

设对象 $\cE$ 有有限 Tor 维数，即 $\cE$ 拟同构于平坦 $\cO_X$-模的有界复形。则一方面 $\cE\otL$ 可延拓到无界导出范畴，另一方面由限制得到从有界导出范畴到自身的函子。我们得到函子
\begin{equation*}
\cE\otL:\mathbf{D}(\cO_X\text{-}\Mod)\longrightarrow\mathbf{D}(\cO_X\text{-}\Mod),\qquad
\cE\otL:\mathbf{D}^{b}(\cO_X\text{-}\Mod)\longrightarrow\mathbf{D}^{b}(\cO_X\text{-}\Mod).
\end{equation*}

注意，若 $\cE$ 属于 $\mathbf{D}^{-}(\cO_X\text{-}\Mod)_{\coh}$（相应地，$\mathbf{D}^{-}(\cO_X\text{-}\Mod)_{\Qcoh}$），则 $\cE\otL$ 把具有（拟）凝聚上同调的对象送到具有（拟）凝聚上同调的对象。

设 $f:X\to Y$ 是诺特概形的态射。直接像函子
\begin{equation*}
f_{*}:\cO_X\text{-}\Mod\longrightarrow\cO_Y\text{-}\Mod
\end{equation*}
是左正合的。由于 $\cO_X$-模的范畴有足够多的内射对象，可知右导出函子
\begin{equation*}
\mathbf{R}f_{*}:\mathbf{D}^{+}(\cO_X\text{-}\Mod)\longrightarrow\mathbf{D}^{+}(\cO_X\text{-}\Mod)
\end{equation*}
存在。此外，此时 $\mathbf{R}f_{*}$ 把子范畴 $\mathbf{D}^{+}(\cO_X\text{-}\Mod)_{\Qcoh}$ 送到子范畴 $\mathbf{D}^{+}(\cO_Y\text{-}\Mod)_{\Qcoh}$。
% [Tutor Note: 页图与文字层此处均把右导出函子的目标范畴印作 $\mathbf{D}^{+}(\cO_X\text{-}\Mod)$；由 $f:X\to Y$ 及紧随其后的“把 $\mathbf{D}^{+}(\cO_X\text{-}\Mod)_{\Qcoh}$ 送到 $\mathbf{D}^{+}(\cO_Y\text{-}\Mod)_{\Qcoh}$”一句，此处应为 $\mathbf{D}^{+}(\cO_Y\text{-}\Mod)$。本译本照录原文，仅注明于此。]

若进一步，$f_{*}$ 有有限的上同调维数，则 $\mathbf{R}f_{*}$ 可延拓到无界复形的范畴。例如，当 $X$ 是有限维诺特概形时就是如此。另一方面，此时（即 $f_{*}$ 有有限上同调维数时），有界导出范畴之间的右导出函子
\begin{equation*}
\mathbf{R}f_{*}:\mathbf{D}^{b}(\cO_X\text{-}\Mod)\longrightarrow\mathbf{D}^{b}(\cO_X\text{-}\Mod)
\end{equation*}
存在。
% [Tutor Note: 同上一处，此式的目标范畴页图亦印作 $\mathbf{D}^{b}(\cO_X\text{-}\Mod)$，应为 $\mathbf{D}^{b}(\cO_Y\text{-}\Mod)$；照录原文。]

要使右导出函子在凝聚层的导出范畴之间有定义，需要对态射添加额外条件。

% [原文 p.529]
\begin{proposition}{1.3.5}
（[16]，III，3.2.1；[17]）设 $f:X\to Y$ 是诺特概形的紧合态射。则函子 $\mathbf{R}f_{*}$ 把子范畴 $\mathbf{D}^{+}(\cO_X\text{-}\Mod)_{\coh}$ 送到子范畴 $\mathbf{D}^{+}(\cO_Y\text{-}\Mod)_{\coh}$。此外，若 $X$ 是有限维的，则对有界导出范畴与无界导出范畴也有类似断言。
\end{proposition}

设 $\cE,\cF\in\mathbf{C}(\cO_X\text{-}\Mod)$ 是 $\cO_X$-模的两个复形。我们由规则
\begin{equation*}
\mathcal{H}om^{n}(\cE,\cF)=\prod_{p}\mathcal{H}om(\cE^{p},\cF^{p+n})
\end{equation*}
定义复形 $\mathcal{H}om^{\cdot}(\cE,\cF)$，其微分为 $d=d_{\cE}+(-1)^{n+1}d_{\cF}$。复形态射之间的同伦可以延拓到局部 $\mathcal{H}om$ 上，于是得到双函子
\begin{equation*}
\mathcal{H}om:\mathbf{H}(\cO_X\text{-}\Mod)^{\mathrm{op}}\times\mathbf{H}(\cO_X\text{-}\Mod)\longrightarrow\mathbf{H}(\cO_X\text{-}\Mod).
\end{equation*}
% [Tutor Note: 原书把层内的 Hom 排作花体 $\mathcal{H}om$，与整体的 $\Hom$ 相区别，其右导出函子记作 $\mathbf{R}\mathcal{H}om$、局部超 Ext 记作 $\mathcal{E}xt$；本译本保留这一花体写法，不用 preamble 中直立体的 \texttt{\textbackslash RHom}、\texttt{\textbackslash Ext}，以免混淆两者。]

由于每个左有界复形都有内射消解，我们得到导出双函子
\begin{equation*}
\mathbf{R}\mathcal{H}om:\mathbf{D}(\cO_X\text{-}\Mod)^{\mathrm{op}}\times\mathbf{D}^{+}(\cO_X\text{-}\Mod)\longrightarrow\mathbf{D}(\cO_X\text{-}\Mod).
\end{equation*}

在此情形下我们定义局部超 Ext
\begin{equation*}
\mathcal{E}xt^{i}(\cE,\cF):=H^{i}(\mathbf{R}\mathcal{H}om(\cE,\cF)).
\end{equation*}

对诺特概形 $X$，若 $\cE$ 与 $\cF$ 是（拟）凝聚的 $\cO_X$-模，则对任一 $i\geqslant 0$，层 $\mathcal{E}xt^{i}(\cE,\cF)$ 也是（拟）凝聚的。

现在若 $\cE\in\mathbf{D}^{-}(\cO_X\text{-}\Mod)_{\coh}$ 且 $\cF\in\mathbf{D}^{+}(\cO_X\text{-}\Mod)_{\coh}$，则 $\mathbf{R}\mathcal{H}om(\cE,\cF)$ 属于 $\mathbf{D}(\cO_X\text{-}\Mod)_{\coh}$。

\yizhu{$\mathcal{E}xt^{i}(\cE,\cF)$ 是 $X$ 上的层，取整体截面才得到群 $\Ext^{i}(\cE,\cF)$，两者不可混用。$\mathcal{H}om$ 只对第一个变元右正合，故 $\mathbf{R}\mathcal{H}om$ 须对第二个变元取内射消解——这解释了为什么它定义在 $\mathbf{D}(\cO_X\text{-}\Mod)^{\mathrm{op}}\times\mathbf{D}^{+}(\cO_X\text{-}\Mod)$ 上，而不是双无界。第 2 章讨论 $\mathbf{D}^{b}(X)$ 上的 $\mathcal{E}xt$ 与 Serre 对偶时，用的正是这种一边无界、一边有界的搭配。【来源：原书 p.529；同 [17]。】}

下面描述本节引入的各导出函子的主要性质与相互关系。考虑两个态射 $f:X\to Y$ 与 $g:Y\to Z$。此时有两个函子 $\mathbf{L}(gf)^{*}$ 与 $\mathbf{L}f^{*}\mathbf{L}g^{*}$，从 $\mathbf{D}^{-}(\cO_Z\text{-}\Mod)$ 到 $\mathbf{D}^{-}(\cO_X\text{-}\Mod)$。于是自然变换
\begin{equation*}
\mathbf{L}(gf)^{*}\xrightarrow{\ \sim\ }\mathbf{L}f^{*}\mathbf{L}g^{*}
\end{equation*}
是同构。这一断言的证明依据如下事实：函子 $g^{*}$ 把平坦 $\cO_Z$-模送到平坦 $\cO_Y$-模（例如见 [17]）。

同样，我们有从 $\mathbf{D}^{+}(\cO_X\text{-}\Mod)$ 到 $\mathbf{D}^{+}(\cO_Z\text{-}\Mod)$ 的函子同构
\begin{equation*}
\mathbf{R}(gf)_{*}\longrightarrow\mathbf{R}g_{*}\mathbf{R}f_{*}.
\end{equation*}
这一断言依据如下事实：$f_{*}$ 把内射层送到 $Y$ 上的松弛层（flabby sheaf），而后者是 $g_{*}$-零调的（见 [17]）。

我们相当频繁地使用的另一些关系，称为投影公式与平坦基变换。

% [原文 p.530]
\begin{proposition}{1.3.6}
（[17]，II.5.6）设 $f:X\to Y$ 是有限维诺特概形之间的态射。则对任意对象 $\cE\in\mathbf{D}^{-}(\cO_Z\text{-}\Mod)$ 与 $\cF\in\mathbf{D}^{-}(\cO_X\text{-}\Mod)_{\Qcoh}$，有函子的自然同构
\begin{equation}
\mathbf{R}f_{*}\cE\otL\cF\xrightarrow{\ \sim\ }\mathbf{R}f_{*}\bigl(\cE\otL\mathbf{L}f^{*}\cF\bigr).
\tag{3}
\end{equation}
\end{proposition}
% [Tutor Note: 命题 1.3.6 中两个对象所属的范畴，页图与文字层一致，均印作 $\cE\in\mathbf{D}^{-}(\cO_Z\text{-}\Mod)$、$\cF\in\mathbf{D}^{-}(\cO_X\text{-}\Mod)_{\Qcoh}$。按式 (3) 的外观（$f:X\to Y$，左端 $\mathbf{R}f_{*}\cE\otL\cF$ 定在 $Y$ 上）此处应为 $\cE\in\mathbf{D}^{-}(\cO_X\text{-}\Mod)$、$\cF\in\mathbf{D}^{-}(\cO_Y\text{-}\Mod)_{\Qcoh}$；原书 $\mathcal O_Z$ 更可能是 $\mathcal O_Y$ 之误。本译本照录原文。]

\begin{proposition}{1.3.7}
（[17]，II.5.12）设 $f:X\to Y$ 是有限维诺特概形之间的有限型态射，$g:Y'\to Y$ 是平坦态射。考虑纤维积方块
\diag{X\times_{Y}Y' \arrow[r,"g'"] \arrow[d,"f'"']
\& X \arrow[d,"f"]\\
Y' \arrow[r,"g"'] \& Y}
此时存在函子的自然同构
\begin{equation}
\mathbf{L}g^{*}\mathbf{R}f_{*}\cE\xrightarrow{\ \sim\ }\mathbf{R}f'_{*}\mathbf{L}g'^{*}\cE,
\tag{4}
\end{equation}
它对任意 $\cE\in\mathbf{D}(\cO_X\text{-}\Mod)_{\Qcoh}$ 成立。
\end{proposition}

我们陈述另一个需要用到关系。

\begin{proposition}{1.3.8}
（[17]，II.5.16）设 $\cE$ 是诺特概形 $X$ 上有限秩局部自由层的有界复形。则下述函子的自然同构
\begin{equation}
\mathbf{R}\mathcal{H}om(\cF,\cG)\otL\cE\xrightarrow{\ \sim\ }\mathbf{R}\mathcal{H}om(\cF,\cG\otL\cE)\xrightarrow{\ \sim\ }\mathbf{R}\mathcal{H}om(\cF\otL\cE^{\vee},\cG)
\tag{5}
\end{equation}
对任意 $\cF\in\mathbf{D}^{-}(\cO_X\text{-}\Mod)$、$\cG\in\mathbf{D}^{+}(\cO_X\text{-}\Mod)$ 成立，其中 $\cE^{\vee}:=\mathbf{R}\mathcal{H}om(\cE,\cO_X)$。
\end{proposition}

\chapterreview{%
第 1 章是全书的地基，它的地位不在结论，而在语言。Orlov 把三样东西在这一章里装配成后面各章通用的工具：Verdier 的三角范畴、经由局部化的导出函子、以及概形上层的导出范畴。第 2 章对凝聚层范畴与 flop 双有理变换的处理、第 3 章“等价必由积上的对象表示”的主定理、以及后文关于 K3 曲面与自等价的讨论，没有一处能绕开本章的词汇表。因此读这一章的正确方式不是记忆，而是随时回查：后面每遇到 $\mathbf{R}f_{*}$、$\otL$、$\mathbf{R}\mathcal{H}om$ 或某个全忠实性断言，其成立条件都写在这里。

知识的脉络是一条单向的收紧链。1.1 从三角范畴的公理 T1--T4 出发，定义正合函子与 Serre 函子，再由局部化与分式演算（L1--L4、引理 1.1.6）把“商范畴”造出来，最后用引理 1.1.10、定义 1.1.11、引理 1.1.12 给出右导出函子在局部化层面的存在判据。1.2 把这套抽象机器落到具体的 $\mathbf{H}(\cA)/\cN$ 上：复形、同伦、映射锥、标准三角，一级一级地把“同调信息”抽出来；随后以内射消解与右适应子范畴把判据变成可计算的形式。1.3 再把 1.2 的结果搬到概形上，沿 $\cO_X$-Mod $\supset$ $\Qcoh(X)$ $\supset$ $\coh(X)$ 三层逐级比较，得到无界、有界、右有界三档等价（命题 1.3.1--1.3.3、推论 1.3.4），然后一次性给出六个基本导出函子与三条运算律。

核心结论可以概括为三点。其一，导出范畴是同伦范畴关于零调复形的局部化，$\mathbf{D}(\cA)=\mathbf{H}(\cA)[\mathrm{Quis}^{-1}]$，而 $\mathbf{D}^{+}$、$\mathbf{D}^{-}$、$\mathbf{D}^{b}$ 分别由消没方向刻画，这一点使“取有界部分”成为范畴层面的操作而非形式约定。其二，三角范畴之间的导出函子能否存在，取决于是否存在适应子范畴：右 $F$-零调对象的存在使 $\mathbf{R}F$ 在 $\mathbf{D}^{+}$ 上处处有定义，而内射对象关于任意加性函子都右适应，这就是“用内射消解算 $\mathbf{R}F$”的全部理由。其三，概形上 $\mathbf{D}^{+}(\Qcoh(X))\to\mathbf{D}^{+}(\cO_X\text{-}\Mod)$ 是全忠实嵌入，像恰为具有拟凝聚上同调的对象；于是可以先在拟凝聚层里做内射消解，再把结论落回凝聚层。$\coh(X)$ 没有足够多内射对象这一缺陷，正是靠这条通道绕过的。

有四处需要提醒。引理 1.2.5 的排版漏了一个 $F$，并把手里的 $\mathbf{H}^{+}(\cR)$ 印成了 $\mathbf{H}^{+}(\cA)$；照字面读会得到“$\mathbf{R}FX$ 同构于一个复形”这种类型不符的断言。第 528 页两处 $\mathbf{R}f_{*}$ 的目标范畴印成 $\cO_X$-Mod，命题 1.3.6 的对象所属范畴又把 $\cO_X$、$\cO_Y$ 写成 $\cO_Z$、$\cO_X$，都应按上下文改正后才好使用。此外，原书 left bounded 指有界在下（属 $\mathbf{D}^{+}$）、right bounded 指有界在上（属 $\mathbf{D}^{-}$），与部分文献相反，读命题 1.3.5 与 1.3.6 时尤须留意。最后，层内的 $\mathcal{H}om$ 与整体的 $\Hom$ 是两回事，$\mathcal{E}xt^{i}$ 是层而非群，这一点在第 2 章讲 $\mathcal{E}xt$ 消没时会被反复用到。

从体例上看，本章是有意“够用即止”的：命题 1.3.1--1.3.3 一律引 [41]（SGA4, Exp. II）与 [17]，不在这里重证；无界情形的一般理论（例如无界复形的内射/平坦消解）只以“有限维”“有限 Tor 维数”“有限上同调维数”这类附加条件一笔带过。这种取舍正是综述的长处：后面每章只调用本章少数几条，无须背负导出范畴的全部技术债。真正耐用的是定义 1.1.3 的 Serre 函子——它在第 2 章化为 $(\,\cdot\,)\otimes\omega_X[n]$，在第 4 章成为 Mukai 格上霍奇等距的来源；以及命题 1.3.6 的投影公式——第 3 章证明 $\Phi_{\cE}$ 全忠实时的每一步换序都靠它。若只允许从本章带走两件东西，就是这两件。%
}
